\section{Governing Equations} \label{cha:Problem_Formulation}
To idealize the problem, we make the following assumptions:
\begin{enumerate}
	\item The displacement is small.
	\item The material is homogeneous and isotropic with modulus of elasticity $E$.
	\item The loading remains within the linear-elastic range, and the material obeys Hooke's law.
	\item The bar is in a state of uniaxial stress. This idealization is not exact because the stress state near the fixed end is
	three-dimensional due to the Poisson effect.\footnote{This ratio is named after Sim\'{e}on Poisson.
	It is defined as the contraction strain, perpendicular to the load, to the extension or axial strain, in the load direction.}
	However, according to Saint-Venant's principle the effects are localized.\footnote{Named after the French elasticity theorist
	Jean Claude Barr\'{e} de Saint-Venant, states that:\quote{The difference between the effects of two different but statically
	equivalent loads becomes very small at sufficiently large distances from load.}}
\end{enumerate}	\text{}
Consider an infinitesimal material element of length $dz$. Enforcing force equilibrium along its longitudinal axis gives
\begin{equation}	\label{eqn:Force_Balance_1}
	-\gls{ABar} \gls{sigma} + \overline{A} \left( z + dz \right) \left[ \gls{sigma} + d \gls{sigma} \right] + \gls{fBar} dz = 0.
\end{equation}
The area for the element's end point may be expressed as
\begin{equation}	\label{eqn:Area_At_The_Right_End}
	\overline{A} \left( z + dz \right) = \gls{ABar} + \frac{ d \overline{A} }{ dz } dz + O \left( dz^2 \right).
\end{equation}
Insert now \cref{eqn:Area_At_The_Right_End} into \cref{eqn:Force_Balance_1} and simplify
\begin{equation}	\label{eqn:Force_Balance_2}
	\gls{ABar} d \sigma + \sigma \frac{ d \overline{A} }{ dz } dz + \frac{ d \overline{A} }{ dz } dz d \sigma + O \left( dz^2
	\right) + \gls{fBar} dz = 0.
\end{equation}
In \cref{eqn:Force_Balance_2}, neglect terms of second order and divide it by $dz$
\begin{equation}	\label{eqn:Force_Balance_3}
		-\frac{ d }{ dz } \Bigl[ \gls{ABar} \sigma \left( z \right) \Bigr] = \gls{fBar}.
\end{equation}
\Cref{eqn:Force_Balance_3} is a differential equation for the stress in the bar. To derive a corresponding equation for the
displacement, consider the following relationships:
\begin{subequations}
	\begin{align}
		\gls{sigma} &= E \gls{epsilon}	\label{eqn:Stress_Strain_1},	\\
		\gls{epsilon} &= \frac{ \gls{uBar} + d \gls{uBar} - \gls{uBar} }{dz} = \frac{ d \overline{u} }{ dz }.
		\label{eqn:Strian_Displacement}
	\end{align}
\end{subequations}
Combining \cref{eqn:Stress_Strain_1,eqn:Strian_Displacement} gives the stress--strain relation
\begin{equation}	\label{eqn:Stress_Strain}
		\gls{sigma} = E \frac{ d \overline{u} }{ dz }.
\end{equation}
A differential equation for the displacement of the bar is then derived by inserting \cref{eqn:Stress_Strain} into
\cref{eqn:Force_Balance_3}
\begin{equation}	\label{eqn:Force_Balance}
		-\frac{ d }{ dz } \biggl[ E \gls{ABar} \frac{ d \overline{u} }{ dz } \biggr] = \gls{fBar}.
\end{equation}
Since the left end of the bar is fixed, its displacement is zero. For a general introduction to FEM, however, we prescribe a
constant displacement at the left boundary. Thus
\begin{equation}
	\gls{uBar} \Big \vert_{z = 0} = \gls{deltaBar}.
\end{equation}
The right end of the bar is subjected to force \gls{PBar}. Thus
\begin{align}
		F \left( z \right) \Big \vert_{z = L} &= \gls{PBar},	\\
		\biggl[ E \gls{ABar} \frac{ d \overline{u} }{ dz } \biggr]_{z = L} &= \gls{PBar}.
\end{align}
\begin{problem} \label{prb:Non_Normalized_Problem}
	The governing equation for the \gls{bar}'s displacement function can be written as
	\begin{subequations} \label{eqn:Non_Normalized_Problem}
		\begin{align}
			-\frac{d}{dz} \Biggl[ \gls{qBar} \frac{d \glssymbol{uBar}}{dz} \Biggr] &= \gls{fBar},
			\quad z \in \left( 0, L \right), \label{eqn:Non_Normalized_Diff_Eq}	\\
			\gls{qBar} &= \glssymbol{E} \gls{ABar},	\label{eqn:Non_Normalized_Rigidity}
		\end{align}
		with boundary conditions
		\begin{align}
			\gls{uBar}\Big\vert_{z=0} &= \gls{deltaBar}, \label{eqn:Non_Normalized_Dirichlet_Cond} \\
			\Biggl[ \gls{qBar} \frac{d \glssymbol{uBar}}{dz} \Biggr]_{z = \gls{L}} &= \gls{PBar}.
				\label{eqn:Non_Normalized_Normal_Cond}
		\end{align}
	\end{subequations}
\end{problem}
%
\begin{remark}
	\Cref{eqn:Non_Normalized_Dirichlet_Cond,eqn:Non_Normalized_Normal_Cond} are both \glspl{boundaryCondition}, but quite
	different in nature. \Cref{eqn:Non_Normalized_Dirichlet_Cond} is a Dirichlet or \gls{essentialCondition} and specifies
	the displacement value on the \gls{leftBoundary}. \Cref{eqn:Non_Normalized_Normal_Cond}, on the other hand, is a normal
	or \gls{naturalCondition}. In a \gls{1DimProblem}, it specifies the boundary flux or traction
	$\overline{q}(L)\,\overline{u}'(L)=\overline{P}$, rather than prescribing the derivative independently of the rigidity.
\end{remark}
%
